% !TeX program = pdflatex
% corpus/docs/torre_fundacao.tex — Quatro fundações ANTES do Corpo de Peano.
% Não se constrói o prédio pelo teto: a teoria (Peano / histerese / música) vem no fim.
% Auto-contido: pdflatex torre_fundacao.tex
\documentclass[11pt,a4paper]{article}
\usepackage[utf8]{inputenc}
\usepackage[T1]{fontenc}
\usepackage[portuguese]{babel}
\usepackage{amsmath,amssymb,amsthm}
\usepackage[margin=2.5cm]{geometry}
\usepackage{booktabs}
\usepackage{xcolor}
\usepackage[hidelinks]{hyperref}
\newtheorem{definicao}{Definição}
\newtheorem{proposicao}{Proposição}
\newtheorem{observacao}{Observação}
\newcommand{\code}[1]{\texttt{\small #1}}
% ─────────────────────────────────────────────────────────────────────────────
%  papers/gkcapa.tex — a capa do Reino, mínima, para os papers em `article`.
%
%  O estilo.tex completo é do `report` (títulos de capítulo, skak, …). Aqui fica
%  só o que a capa precisa, para o paper continuar a compilar sozinho com
%  pdflatex. O tradutor próprio (tests/tex) carrega o estilo.tex da raiz / o
%  symlink papers/estilo.tex e avalia o mesmo `\gkcapa`.
% ─────────────────────────────────────────────────────────────────────────────
\definecolor{ouro}{HTML}{B8912F}
\definecolor{ouroclaro}{HTML}{F3E9CE}
\definecolor{tinta}{HTML}{1E1B16}
\definecolor{regua}{HTML}{6E6858}
\providecommand{\gknota}{\fontsize{7.62}{11.05}\selectfont}
\providecommand{\gktexto}{\fontsize{10.50}{15.22}\selectfont}
\providecommand{\gksub}{\fontsize{12.33}{17.87}\selectfont}
\providecommand{\gksec}{\fontsize{14.47}{20.98}\selectfont}
\providecommand{\gkcap}{\fontsize{16.99}{24.63}\selectfont}
\providecommand{\gktit}{\fontsize{23.42}{33.95}\selectfont}
\providecommand{\Prj}{\mathbb{P}}
\providecommand{\Trf}{\mathcal{T}}
\providecommand{\gkcapa}[3]{%
\title{\vspace{-0.6cm}
{\color{ouro}\rule{\textwidth}{1.4pt}}\\[6mm]
\textbf{\gktit\color{tinta} Reino Dourado}\\[3mm]{\gksec\itshape\color{regua} Golden Kingdom}\\[8mm]
{\gkcap\scshape\color{ouro} #1}\\[5mm]
{\gksub\color{regua} #2}\\[2mm]
{\gktexto\color{regua}\itshape #3}\\[7mm]
{\color{ouro}\rule{\textwidth}{0.6pt}}\\[9mm]{\gknota\color{regua}\begin{minipage}{\textwidth}\setlength{\parindent}{0pt}%
\textbf{Aviso de Isenção Algorítmica:} O universo, as leis e as entidades contidas nestas páginas ---
incluindo felinos matriciais, aves de rapina algébricas e amuletos de controle de fluxo --- são
construções de pura ficção formal. Esta obra não descreve a realidade física, não serve como
engenharia reversa do cosmos e não constitui doutrina religiosa. Qualquer semelhança com bugs reais do seu
sistema operacional é mera coincidência (ou falta de reza). Prossiga por sua própria conta e
risco.\end{minipage}}}
\author{}\date{}}

\begin{document}
\gkcapa{Torre de Fundação}
       {Corpos $\to$ Dinâmica $\to$ Topologia $\to$ Hilbert/Clifford}
       {antes do Corpo de Peano --- a teoria chega por último}
\maketitle

\begin{abstract}
\noindent Não se constrói um prédio pelo teto. Este documento é a
\textbf{torre abaixo}: quatro fundações, cada uma a responder a uma
pergunta que a torre precisa resolver. Só depois se entra no
\emph{Corpo de Peano} (\code{papers/corpo\_topologico.tex}), no Teorema
Central Gentil\,$\leftrightarrow$\,Hurwitz, na música e na histerese.
\[
\boxed{
\text{Corpos}
\;\to\;
\text{Dinâmica}
\;\to\;
\text{Topologia}
\;\to\;
\text{Hilbert/Clifford}
\;\to\;
\text{Corpo de Peano}
\;\to\;
\text{Teorema Central}.
}
\]
Não é um capítulo universitário genérico: é a escada até a
arquitectura aparecer como consequência.
\end{abstract}

\medskip
\noindent\textbf{Escada de perguntas.}
\begin{center}\small
\begin{tabular}{@{}lp{8.6cm}@{}}
\toprule
Fundação & Pergunta \\
\midrule
Corpos & O que pode ser composto? \\
Dinâmica & Como isso evolui? \\
Topologia & Onde isso pode ir sem perder continuidade? \\
Hilbert/Clifford & Como medir e orientar o movimento? \\
Peano (depois) & Como tudo isso se realiza numa torre? \\
Música (depois) & Como a torre pode ser lida? \\
Teorema Central (fim) & Por que as duas leituras são a mesma estrutura? \\
\bottomrule
\end{tabular}
\end{center}

%----------------------------------------------------------------------
\section{Fundação I --- Teoria de corpos}\label{sec:fund-corpos}
\subsection*{Pergunta: o que pode ser composto?}

Antes de Maestro ou partitura, o leitor precisa de saber o que é um
\textbf{corpo}: um sítio onde se soma, multiplica, e onde os não-nulos
têm inverso.

\begin{definicao}[corpo --- mínimo operacional]
Um conjunto $K$ com $+$ e $\times$ é um \textbf{corpo} quando $(K,+)$ é
grupo abeliano, $(K\setminus\{0\},\times)$ é grupo abeliano, e a
distribuição liga as duas operações. O zero é o neutro de $+$; o um é
o neutro de $\times$.
\end{definicao}

\noindent O que a torre vai precisar já aqui:
\begin{itemize}
\item \textbf{operações e inversores} --- $a\mapsto -a$, $a\mapsto a^{-1}$
quando $a\neq0$;
\item \textbf{extensão e redução} --- passar de $K$ a $L\supset K$ (ou
projectar de volta) sem inventar axiomas novos;
\item \textbf{divisão / Hurwitz} --- a família de álgebras com norma
multiplicativa (reais, complexos, quaterniões, octoniões) é o teto
clássico da composição; o Corpo de Peano \emph{realizará} essa escada
em andares $d_{k+1}=2d_k$, mas só depois.
\end{itemize}

\begin{observacao}
Aqui o leitor aprende \textbf{o que é o corpo} antes de descobrir que o
sistema o realiza musicalmente. Sem esta fundação, ``corpo de Peano''
soa a metáfora.
\end{observacao}

%----------------------------------------------------------------------
\section{Fundação II --- Dinâmica de sistemas}\label{sec:fund-dinamica}
\subsection*{Pergunta: como isso evolui?}

Composto o objecto, pergunta-se o que acontece quando se move:
\[
\boxed{
X_{k+1}=F(X_k).
}
\]

\noindent Entram naturalmente:
\begin{itemize}
\item \textbf{órbitas} --- $\{X_0,X_1,\ldots\}$ sob $F$;
\item \textbf{estados} --- o ponto $X_k$ no espaço de fase;
\item \textbf{reversibilidade} --- existência de $F^{-1}$ (ou involução
local): a volta pode fechar;
\item \textbf{estabilidade} --- órbitas que permanecem perto;
\item \textbf{expansão / contração} --- taxas $\lambda^{+}$, $\lambda^{-}$
(Lyapunov); numa volta fechada dualizada,
$\lambda^{+}+\lambda^{-}=0$;
\item \textbf{memória} --- o estado seguinte depende do caminho, não só
da posição instantânea (histerese clássica: entrada $\neq$ saída).
\end{itemize}

\begin{proposicao}[pré-imagem operacional]
Maestro, Metrónomo e Controle de Histerese \textbf{não} se inventam
aqui: são nomes que a torre dará, mais tarde, a realizações desta
dinâmica (projectar / ler / seleccionar). Nesta fundação basta:
estado, órbita, volta, expansão/contração.
\end{proposicao}

%----------------------------------------------------------------------
\section{Fundação III --- Topologia}\label{sec:fund-topologia}
\subsection*{Pergunta: onde isso pode ir sem perder continuidade?}

A dinâmica precisa de \textbf{sítio}. A topologia responde com
vizinhança, continuidade e borda:
\[
\boxed{
\mathcal{N}(x),\qquad
\partial B,\qquad
\pi\colon A\to A'
\text{ (projecção / retração)}.
}
\]

\noindent Objectos mínimos:
\begin{itemize}
\item \textbf{vizinhança} $\mathcal{N}(x)$ --- o que ainda conta como
``perto'' de $x$;
\item \textbf{continuidade} --- $F$ leva vizinhos a vizinhos;
\item \textbf{borda} $\partial B$ --- interface entre dentro e fora
(zona cinzenta da histerese clássica);
\item \textbf{retração / projecção} --- $\pi\circ\iota=\mathrm{id}$
quando a torre descer andares ($\pi_k\colon A_{k+1}\to A_k$ virá no
Peano);
\item \textbf{identificação} --- colar pontos sem destruir a estrutura
local.
\end{itemize}

\begin{observacao}
Com topologia, a histerese deixa de parecer um mecanismo inventado para
o corpus: passa a ser \textbf{consequência natural} de uma borda com
direcção de passagem. O Controle de Histerese escolherá variantes
dentro de $\mathcal{V}(B)$, mas isso é o teto --- ainda não.
\end{observacao}

%----------------------------------------------------------------------
\section{Fundação IV --- Hilbert / Clifford}\label{sec:fund-hilbert}
\subsection*{Pergunta: como medir e orientar o movimento?}

Falta a \textbf{régua} e a \textbf{orientação}:
\[
\boxed{
\langle x,y\rangle,\qquad
|x|=\sqrt{\langle x,x\rangle},\qquad
xy+yx=2\langle x,y\rangle.
}
\]

\noindent Hilbert dá o produto interno e a norma; Clifford dá a álgebra
que \textbf{orienta} o movimento (reflexão, inversão, conservação):
\begin{itemize}
\item \textbf{norma} --- comprimento do passo;
\item \textbf{dualidade} --- formas lineares / covectores;
\item \textbf{reflexão} --- a volta métrica (espelho);
\item \textbf{orientação} --- sentido do MOVE;
\item \textbf{inversão} --- $x\mapsto x^{-1}$ quando a norma permite;
\item \textbf{conservação} --- $|xy|=|x||y|$ nas álgebras de Hurwitz.
\end{itemize}

\begin{proposicao}[ponte]
Esta fundação prepara directamente Gentil\,$\leftrightarrow$\,Hurwitz:
contar e integrar tornam-se duais \emph{pela medida} --- mas o Teorema
Central só se enuncia \textbf{depois} do Corpo de Peano ter a torre
$A_k$ e as retrações $\pi_k$.
\end{proposicao}

%----------------------------------------------------------------------
\section{O que ainda não se faz aqui}\label{sec:fund-nao}

Não se enuncia ainda:
\begin{itemize}
\item o Teorema do Maestro / Metrónomo;
\item o Teorema da Histerese nem o do Controle;
\item a partitura como linguagem de unificação;
\item Gentil\,$\leftrightarrow$\,Hurwitz como teorema fechado.
\end{itemize}
Esses objectos vivem em \code{papers/corpo\_topologico.tex} e só fazem
sentido \textbf{depois} destas quatro fundações. A ordem operacional
do projecto espelha a mesma disciplina no corpus da assistente:
primeiro falar (banal); depois subir complexidade; a doutrina no fim
(\code{corpus/fala/conversa.tex}, \code{tools/vizinhanca.sh}).

\section{Passagem ao Corpo de Peano}

Quando as quatro perguntas tiverem resposta mínima, pode dizer-se:
\[
\boxed{
\text{Agora temos os objectos necessários para construir o Corpo de Peano.}
}
\]
A torre começa:
\[
A_0
\;\xleftarrow{\pi_0}\;
A_1
\;\xleftarrow{\pi_1}\;
A_2
\;\xleftarrow{\pi_2}\;
\cdots
\qquad
d_{k+1}=2d_k.
\]
Aí o Maestro é a realização dinâmica das retrações; o Metrónomo a
leitura métrica; o Controle a selecção na vizinhança. Só no fim o
Teorema Central:
\[
\boxed{\text{Gentil}\;\longleftrightarrow\;\text{Hurwitz}}
\]
e a música deixa de ser metáfora --- passa a ser realização gráfica de
uma estrutura que o leitor já encontrou na fundação.

\bigskip
\noindent\textbf{Leitura.} Este paper $\to$
\code{papers/corpo\_topologico.tex} $\to$
\code{corpus/docs/partitura.tex} / \code{corpus/docs/dualsort.tex}.
Corpus (ainda sem aula): \code{tools/bench\_vizinhanca.sh}.

\end{document}
